%% it/sections/03_fazioni.tex
%% ============================================================
%% 03_fazioni.tex — Le quattro fazioni / The four factions
%% ============================================================

\chapter{Le Fazioni / The Factions}
\markboth{Le Fazioni}{The Factions}

\lettrine[lines=3]{Q}{uattro} fazioni si muovono nella Palermo del febbraio
1282, ignare o consapevoli l'una dell'altra, legate da interessi che si
sovrappongono e si contraddicono. Nessuna controlla completamente la
situazione. Tutte, in modi diversi, spingeranno il 30 marzo verso
ciò che diventerà.

\nota{Per il GM Coordinatore}{
  Le fazioni non si conoscono necessariamente tra loro. Parte della
  tensione narrativa nasce proprio da questa opacità: i Congiurati
  credono di controllare i tempi, il Clero crede di gestire i simboli,
  i Baroni credono di negoziare il futuro, il Popolo non crede in niente
  — ma agisce. Le interferenze tra i tavoli devono essere gestite con
  delicatezza: usate i Momenti di Congiunzione per far filtrare le
  conseguenze, non le cause.
}

\secsep

%% ---- FAZIONE 1: CONGIURATI ----
\section{I Congiurati}
\textcolor{oro}{\cinzel\large La Rete di Giovanni da Procida}

\filetto

\textbf{Punto di ingresso:} In medias res — i personaggi sono già dentro.

\textbf{Obiettivo di sessione:} Assicurarsi che la rivolta scoppi prima
che la flotta di Carlo salpi. Il piano prevede un segnale coordinato,
ma i dettagli sono nei dettagli — e i dettagli stanno cedendo.

\textbf{Tema narrativo:} La fiducia in una causa non protegge dal
tradimento. Chi comanda da lontano (da Procida, Pietro d'Aragona, Michele
Paleologo) ha interessi che non coincidono con quelli di chi rischia
la vita a Palermo.

\subsection{Struttura della Fazione}

I Congiurati sono una rete, non un'organizzazione. I personaggi si
conoscono tra loro ma non conoscono necessariamente tutti i nodi
della rete. Sanno che esiste un piano, sanno il loro pezzo — ma
non vedono il quadro completo. Questa opacità è intenzionale e
narrativamente produttiva.

\subsection{Punteggio di Fazione}

\begin{itemize}[nosep]
  \item \textbf{+1} per ogni contatto baronale attivato e affidabile
  \item \textbf{+1} se il canale con la Sicilia orientale (Messina) è aperto
  \item \textbf{+1} se l'identità della rete non è stata compromessa
  \item \textbf{+1} se il segnale coordinato è stato definito e comunicato
  \item \textbf{-1} per ogni membro della rete catturato o ucciso
  \item \textbf{-1} se le autorità angioine sono in allerta specificamente
    per la cospirazione (non generica)
\end{itemize}

\nota{Per il GM dei Congiurati}{
  I personaggi di questo tavolo sanno più degli altri su cosa sta per
  succedere — ma questa conoscenza è un peso, non un vantaggio. Sono
  responsabili di ciò che accadrà. Se la rivolta scoppia nel modo sbagliato,
  è colpa loro. Se non scoppia, è colpa loro. La pressione psicologica
  è parte del design.
}

\secsep

%% ---- FAZIONE 2: BARONI ----
\section{I Baroni Siciliani}
\textcolor{oro}{\cinzel\large La Nobiltà tra Due Fuochi}

\filetto

\textbf{Punto di ingresso:} Reclutati — all'inizio della sessione vengono
avvicinati da un emissario (che può essere un Congiurato, un frate,
o un misterioso intermediario) con una proposta che non possono ignorare.

\textbf{Obiettivo di sessione:} Decidere da che parte stare — e sopravvivere
alla scelta. I baroni siciliani hanno tutto da perdere: le loro terre, i
loro titoli, le loro famiglie. Hanno tutto da guadagnare: la possibilità
di liberarsi dal giogo angioino. Il problema è che entrambe le opzioni
sono pericolose.

\textbf{Tema narrativo:} Il privilegio non è protezione. Avere potere
significa avere responsabilità verso chi dipende da te — e queste
responsabilità si contraddicono a vicenda.

\subsection{La Posizione dei Baroni}

La nobiltà siciliana del 1282 era in una posizione strutturalmente
instabile. Molti di loro avevano collaborato con gli Angioini per
preservare le proprie terre; alcuni avevano ottenuto nuovi privilegi;
pochi si erano apertamente opposti. Ma nessuno era soddisfatto.
Il sistema fiscale angioino colpiva le loro rendite; i funzionari
francesi ignoravano i loro diritti tradizionali; le requisizioni
militari svuotavano i loro vassalli.

Alaimo da Lentini, il più importante barone filo-ribelle della Sicilia
orientale, è il loro punto di riferimento naturale — ma è lontano, e
la sua posizione è ambigua quanto la loro.

\subsection{Punteggio di Fazione}

\begin{itemize}[nosep]
  \item \textbf{+1} se almeno tre casate baronali hanno preso posizione
  \item \textbf{+1} se i propri vassalli sono stati preparati e
    stanno tenendo
  \item \textbf{+1} se si è stabilito un canale sicuro con i Congiurati
    o con il Clero
  \item \textbf{+1} se le proprie famiglie sono al sicuro
  \item \textbf{-1} per ogni casata baronale passata dalla parte angioina
  \item \textbf{-1} se un personaggio è stato catturato o compromesso
\end{itemize}

\secsep

%% ---- FAZIONE 3: CLERO ----
\section{Il Clero}
\textcolor{oro}{\cinzel\large I Canali della Chiesa}

\filetto

\textbf{Punto di ingresso:} Misto — alcuni personaggi sono già consapevoli
di ciò che si prepara (un vescovo informato, un frate usato come
corriere), altri vengono trascinati dagli eventi.

\textbf{Obiettivo di sessione:} Gestire il momento senza che la Chiesa
ne esca compromessa — e allo stesso tempo assicurarsi che ciò che è
giusto accada. La tensione tra fedeltà istituzionale e coscienza
individuale è il cuore di questo tavolo.

\textbf{Tema narrativo:} L'obbedienza non assolve dalla responsabilità.
Servire Dio e servire Roma non è la stessa cosa. E servire il proprio
popolo non è la stessa cosa di servire nessuno dei due.

\subsection{Il Clero come Rete}

Le chiese e i conventi di Palermo sono nodi di comunicazione naturali.
I frati circolano liberamente; le confessioni sono inviolabili; i
campanili sono i punti più alti della città. Il Clero non ha eserciti,
ma ha qualcosa di più prezioso: può muovere simboli, e i simboli,
in questo momento storico, muovono le persone.

La chiesa di Santo Spirito — fuori dalle mura, teatro della
processione di Pasqua — è il centro fisico della sessione per questo tavolo.

\subsection{Punteggio di Fazione}

\begin{itemize}[nosep]
  \item \textbf{+1} se la processione di Santo Spirito si svolge
    come previsto (massima concentrazione di persone)
  \item \textbf{+1} se i canali di comunicazione tra i quartieri
    sono attivi e sicuri
  \item \textbf{+1} se si è riusciti a tenere la Chiesa formalmente
    fuori dalla cospirazione (plausible deniability)
  \item \textbf{+1} se almeno un alto prelato è pronto a legittimare
    la rivolta a posteriori
  \item \textbf{-1} se un membro del clero è stato arrestato o
    torturato dai francesi
  \item \textbf{-1} se le autorità angioine hanno cancellato o
    militarizzato la processione
\end{itemize}

\nota{Per il GM del Clero}{
  Questo tavolo ha accesso ai simboli. I personaggi possono fare cose
  che nessun altro può fare: suonare le campane al momento sbagliato
  (o giusto), aprire o chiudere le porte di una chiesa, confessare
  e assolvere. Usate questa asimmetria. Se è attivo il layer Passionale,
  il Clero dispone anche dell'azione \textit{Legare e Sciogliere} (vedi
  «La Passione dei Vespri»): può forgiare o spezzare le Connessioni
  altrui spendendo Passione.
}

\secsep

%% ---- FAZIONE 4: POPOLO ----
\section{Il Popolo}
\textcolor{oro}{\cinzel\large La Miccia}

\filetto

\textbf{Punto di ingresso:} Trascinati — i personaggi non sanno nulla
di nessuna cospirazione. Sono artigiani, pescatori, donne del mercato,
giovani senza prospettive. Vengono reclutati o coinvolti da qualcuno
che si fida di loro — o che ha bisogno di loro.

\textbf{Obiettivo di sessione:} Sopravvivere. Proteggere chi si ama.
E forse, nel processo, capire che quello che sta per accadere non è
solo una notte di sangue — è qualcosa che cambierà tutto.

\textbf{Tema narrativo:} La storia la fanno le persone normali, non
gli eroi. Il Vespro non è scoppiato perché qualcuno aveva un piano:
è scoppiato perché la gente comune aveva raggiunto il limite. Questo
tavolo gioca quel limite.

\subsection{La Vita Quotidiana sotto gli Angioini}

I personaggi del Popolo vivono la pressione angioina in modo concreto
e quotidiano: tasse impossibili, requisizioni, soldati che perquisiscono
le donne, funzionari che prendono quello che vogliono. Non hanno
accesso alle grandi manovre politiche — ma sono esattamente quelli
che saranno in piazza la sera del 30 marzo.

\subsection{Punteggio di Fazione}

\begin{itemize}[nosep]
  \item \textbf{+1} se i quartieri popolari sono informati e mobilitati
    (anche inconsapevolmente)
  \item \textbf{+1} se un sopruso angioino pubblico è stato documentato
    o testimoniato in modo da poter essere usato come detonatore
  \item \textbf{+1} se le reti di vicinato (i \textit{hara}) sono
    pronte ad agire in modo coordinato
  \item \textbf{+1} se i personaggi sono riusciti a proteggere
    qualcuno di vulnerabile (una famiglia, un anziano, un bambino)
    dalle violenze preliminari
  \item \textbf{-1} se una figura chiave della comunità è stata
    eliminata o intimidita dai francesi
  \item \textbf{-1} se i quartieri sono divisi o in conflitto tra loro
\end{itemize}

\nota{Per il GM del Popolo}{
  Questo è il tavolo emotivamente più intenso. I personaggi non hanno
  risorse, non hanno piani, non hanno protezioni. Hanno solo l'un l'altro.
  Usate i piccoli dettagli: il nome di un vicino, il profumo del pane,
  il suono di una perquisizione notturna. La storia grande entra
  dalla porta di servizio.
}

\filettodoppio

%% VERSIONE INGLESE
